\documentclass[11pt,a4paper]{article}
\usepackage[T1]{fontenc}
\usepackage[utf8]{inputenc}
\usepackage{amsmath,amssymb,xcolor}
\usepackage[margin=25mm]{geometry}
\usepackage[colorlinks=true,urlcolor=blue]{hyperref}
\definecolor{nexorablue}{HTML}{2459B5}
\setlength{\parindent}{0pt}
\setlength{\parskip}{7pt}
\setlength{\emergencystretch}{2em}
\renewcommand{\familydefault}{\sfdefault}
\begin{document}

{\small\bfseries\color{nexorablue} NEXORA CPGE}\par
{\small MP / NEX-STR-06 / Corrigé}\par
{\LARGE\bfseries Unités d’une algèbre triangulaire simple}\par
\vspace{0.25em}\hrule\vspace{0.6em}

\section*{Stabilité}
L’ensemble est le sous-espace engendré par $I$ et $N$. Le produit $(aI+bN)(cI+dN)=acI+(ad+bc)N$ appartient encore à cet ensemble. Comme $I$ y appartient, il s’agit d’une sous-algèbre unitaire.

\section*{Unités}
Un inverse exige $ac=1$, donc $a\ne0$. Dans ce cas $c=1/a$ et $d=-b/a^2$ annulent le coefficient de $N$. Le calcul donne l’inverse dans les deux ordres.


\end{document}
